\chapter{Einleitung}


\section{Ausgangslage und Problemstellung}

\subsection{Situation}
In den letzten Jahren wurde der Umgang mit künstlicher Intelligenz von einem "Kann" zu einem "Muss" und
hat so zu einer fundamentalen Veränderung in der IT-Landschaft geführt.
Gerade LLM Models haben als transformative Technologie die Art und Weise, wie wir Arbeiten,
wie wir Informationen beschaffen und wie wir Probleme lösen, revolutioniert. Sie wird nicht nur als Beschleuniger in der Softwareentwicklung
sondern auch als Effizienzsteigerer in der Kundenkommunikation genutzt \footcite{Goodfellow-et-al-2016}.
Besonders beliebt ist dabei die Integration von bestehenden LLM-APIs in eigene Softwarelösungen, da Unternehmen so schnell und kostengünstig
von den Vorteilen der künstlichen Intelligenz profitieren können. \footcite{openai2024gpt4technicalreport}


% Zur verdeutlichung ein gängiges Beispiel: Ein IT-Administrator analysiert regelmäßig die Logdateien eines Servers,
% welcher einen kritischen Fehler bei der Verbindung mit einer Datenbank meldet. Um das Problem schneller zu lösen,
% kopiert er die Logdateien und verwendet sendet sie an eine externe LLM API. Dabei achtet er nicht darauf,
% kritische unternehmensinterne Daten, wie z.B. IP-Adressen, Servernamen oder Passwörter zu entfernen.
% Ohne Data-Masking werden diese Daten ungeschützt an den API Provider gesendet, der sie für Trainingszwecke,
% oder sogar als Sicherheits-insights in das Unternehmen nutzen kann.


\subsection{Task}
Es ist ersichtbar, dass die Integration von LLM-APIs in Softwarelösungen zwar starke Effizienzsteigerungen ermöglicht,
aber auch ein oft übersehenes Risiko mit sich bringt. Das Übetragen von kritischen Unternehmensdaten an Dritte, korreliert unmittelbar
mit dessen wirtschaftlichen Erfolg. TODO QUELLE Wichtig dabei ist zu verstehen, dass die Verschlüsselung der Anfragen über TLS zwar die
Integrität und Vertraulichkeit der Daten sichert, nicht aber den Zugriff des Anbieters auf die Inhalte des Unternehmens. \footcite{europaDeployingPseudonymisation}
Das bedeutet, dass die API Provider die Daten auch als Analysezwecke des Unternehmens nutzen können was regulatorische
und wettbewerbswirksame Implikationen haben kann.

\subsection{Action}

\subsection{Result}
